\section{Signed data, integral Fermat stages, and additive-polynomial lifts}

We distinguish three objects: the action of actual profinite signed
data; the integral group schemes obtained by extending abelian
lattices in finite extensions to vector groups; and arbitrary
automorphisms of those group schemes after reduction modulo $2$.
The last can contain nonlinear additive polynomials which are not
induced by the first. This distinction is essential for task H.3.

\begin{statement}{H.1}{Open}{Images of the retained normalizer and conjugator data}
The complete images of $N_\rho^*$, $N_\sigma^*$, $G_D$, and the
two fibre products on the integral $q_2$ stage and every higher
relative Fermat stage are not determined. H.4 verifies the finite
scheme-theoretic image which every compatible $q_2$ lift must
preserve. H.7 gives typed action formulas and exact tests for a
\emph{given} retained datum, and H.8 excludes the nonlinear
infinitesimal polynomial direction from all actual profinite source
images. These do not identify the requested source images with the
full ambient lift functor.
\end{statement}

\begin{statement}{H.2}{Open}{Stagewise source kernels}
The kernels inside each of the five retained source functors are not
given by an explicit presentation at every stage. For the underlying
automorphism action, H.7 proves the exact matrix-and-cocycle kernel
test and determines its joint kernel on all full profinite abelian
Fermat stages: it is the kernel of the action on the metabelian
quotient. Pulling these tests back to a retained-data source is
valid, but does not by itself compute the fibres of that source's
map to automorphisms or its individual stagewise kernels.
\end{statement}

\begin{statement}{H.3}{Open}{Extension of every additive-polynomial lift}
Every coordinatewise additive-polynomial automorphism fixing $1$
extends compatibly through all integral \emph{abelian-kernel}
Fermat stages in characteristic $2$, by H.5. This is an unbounded
positive result within that category.

Extension through all higher nonabelian relative stages is not
proved. H.6 gives a sharp obstruction: the particular lift
$f(T)=T+\epsilon(T^2-T)$ cannot act coordinatewise on the
abelianization of a Heisenberg quotient with its ordinary
commutator pairing. To apply this obstruction to the full requested
tower, one must construct its specified integral nonabelian stage
and show that it has that compatible quotient. A characteristic-zero
relative completion does not by itself supply such an integral
model at $2$, and none is identified here. Thus H.6 is not labelled
a counterexample to the entire task H.3.
\end{statement}

\begin{statement}{H.4}{Known}{The integral $q_2$ obstruction rechecked}
Let $C=\mathrm{SL}_2(\mathbf Z)$, with
\[
 \rho=\begin{pmatrix}0&1\\-1&1\end{pmatrix},\qquad
 \sigma=\begin{pmatrix}0&-1\\1&0\end{pmatrix},\qquad
 \omega=-I,
\]
and let $x=\omega(\sigma\rho)^2$, $y=\omega(\rho\sigma)^2$.
Then $F=\langle x,y\rangle$ is free of rank two and index twelve,
and $K=\ker(F\to(\mathbf Z/2)^2)=\Gamma(4)$ has rank five.
In the Schreier basis
\[
 x^2,\quad yxy^{-1}x^{-1},\quad y^2,\quad
 xyxy^{-1},\quad xy^2x^{-1},
\]
the lattice map $K^{\mathrm{ab}}\to F^{\mathrm{ab}}$ is
\[
       J=\begin{pmatrix}2&0&0&2&0\\0&0&2&0&2\end{pmatrix}.
\]
For the associated integral cocycle group schemes
$\mathscr C_2\to\mathscr C_0$, reduction modulo $2$ kills the
entire rank-five vector radical. Its scheme-theoretic image is the
constant subgroup $H_{48}=C/K$, with equations
$X^2-X=Y^2-Y=0$ on every component of $\mathscr C_0$.

Consequently the minimal-stage transformation
$u_{t,e}(\rho)=(t,t+e)\rho$,
$u_{t,e}(\sigma)=(t,t)\omega^e\sigma$ can lift only if
$t^2=t$. For $A=\mathbf F_2[\epsilon]/\epsilon^2$,
$t=\epsilon$, $e=0$ is an actual scheme-theoretic obstruction.
Conversely the locus $t^2=t$ lifts by the indicated discrete
conjugation and reversing automorphisms, decomposing the base by
the idempotents if necessary.
\end{statement}
\begin{proof}[Independent checks of the supplied result]
The modular amalgam $C_6*_{C_2}C_4$ and direct matrix multiplication
give the quotient of order twelve after killing $x,y$; its two
finite factors inject. The kernel acts freely on the amalgam tree,
so is free, and its Euler characteristic gives rank two. Reduction
modulo four gives $48$ matrices, with the images of $x,y$ a Klein
four subgroup. Thus the stated index-four kernel in $F$ is
$\Gamma(4)$, and Schreier's formula gives rank five. Rewriting with
the four cosets $1,x,y,xy$ gives the displayed basis, whose exponent
sums give $J$ without a rational rank calculation.

For a lattice extension of a finite group $H$, choose word
representatives $s_h$, write $A_h$ for conjugation on the lattice,
and $c(h,k)$ for $s_hs_ks_{hk}^{-1}$ in its abelianization. The
group scheme is the disjoint union of copies of the vector group,
with law
\[
 (v,h)(w,k)=(v+A_hw+c(h,k),hk).
\]
The projection has form $(v,h)\mapsto(Jv+m_h,p(h))$ with integral
$m_h$. Modulo $2$ the vector term vanishes; the kernel of the
discrete map $C\to\mathscr C_0(\mathbf F_2)$ is exactly $K$.
Thus its image is the stated constant $48$-element subgroup, the
four-point grid in each of twelve components. The projection from
the $48$ vector components is faithfully flat, so this is also the
scheme-theoretic image after every base change.

A compatible automorphism must preserve that image. Applying it
to $\rho$ forces $t^2-t=0$. At $t=\epsilon$ this equation has
value $\epsilon\ne0$, including over a faithfully flat extension
of $A$. For idempotent $t$, the lifts are obtained by the
conjugation by $x$ on the $t=1$ component and the identity on
the $t=0$ component, together with the reversing involution when
$e=1$.

The attached stand-alone program was copied with isolated output
names and executed as \texttt{code/fgh\_q2\_recheck.py}. It checks
all $48^3=110592$ integral cocycle identities, $48^2=2304$
projection identities, the six boundary vectors, and the mod-$2$
linear equations. Over the dual numbers its finite enumeration
finds $256$ fine parameter points, $8$ minimal automorphisms,
and $4$ liftable minimal automorphisms, hence $128$ liftable
parameter points. These numbers concern that explicitly specified
minimal functor, not the full retained-data sources of H.1.
\end{proof}

\begin{statement}{H.5}{Proved}{All abelian Fermat stages admit the same nonlinear lift}
For $n\geq1$ let $K_n=\ker(F\to(\mathbf Z/n)^2)$ and form the
integral cocycle group scheme $\mathscr C_n$ from
$C/[K_n,K_n]$. It has finite component group $C/K_n$ and vector
radical of rank $n^2+1$. For an $\mathbf F_2$-algebra $R$, let
$f\in\operatorname{Aut}_{R\text{-groups}}(\mathbf G_{a,R})$
be an additive polynomial with $f(1)=1$. Then
\[
       \Phi_{n,f}(v,h)=(f(v_1),\ldots,f(v_{n^2+1}),h)
\]
is a group-scheme automorphism. These automorphisms commute with
every tower projection $\mathscr C_m\to\mathscr C_n$ for $n\mid m$,
with change of integral Schreier coordinates, and with every
integral boundary-vector subgroup. They fix every point coming
from a discrete word of $C$.
\end{statement}
\begin{proof}
All matrices $A_h$, cocycle vectors $c(h,k)$, and affine projection
coefficients are integers, reducing to $\mathbf F_2$ in $R$.
Additivity and $f(1)=1$ imply
\[
 f(A_hv)=A_hf(v),\qquad f(c(h,k))=c(h,k).
\]
Substitution in the cocycle law proves that $\Phi_{n,f}$ is a
homomorphism. Its inverse is coordinatewise $f^{-1}$, which also
fixes $1$. A projection is $(v,h)\mapsto(Jv+m_h,p(h))$ with
integral $J,m_h$; the same identities prove compatibility.
Changing integral lattice bases and word representatives is another
such affine transformation, so the construction is independent of
these choices. A boundary vector $p\in\mathbf Z^{n^2+1}$ satisfies
$f(tp)=f(t)p$, proving preservation of its vector subgroup. Finally
all vector coordinates of a discrete word are integers, which
reduce to $0$ or $1$ and are fixed by $f$.

In particular, over $\mathbf F_2[\epsilon]/\epsilon^2$,
\[
             f(T)=T+\epsilon(T^2-T)
\]
is additive, fixes $0,1$, and is its own inverse. The term
$\epsilon T^2$ is nonzero as a polynomial. Although this polynomial
is invisible on the four points of the dual-number ring itself,
it is not the identity morphism of group schemes: it is detected
after adjoining an element of $\mathbf F_4$. The executed
\texttt{code/fgh\_additive\_lifts.py} verifies all $256$ additive
identities on $\mathbf F_4[\epsilon]/\epsilon^2$ and its
compatibility with all $48$ actual $q_2$ action matrices and all
$2304$ cocycle values. The proof above, rather than a finite-stage
induction, establishes compatibility for every $n$.
\end{proof}

\begin{statement}{H.6}{Proved}{A class-two obstruction to nonlinear additive lifts}
Let $R=\mathbf F_4[\epsilon]/\epsilon^2$, where
$\omega^2+\omega+1=0$, and let $\mathscr H_R$ be the Heisenberg
group scheme with multiplication
\[
 (a,b,c)(a',b',c')=(a+a',b+b',c+c'+ab').
\]
There is no group-scheme automorphism of $\mathscr H_R$ inducing
$(a,b)\mapsto(f(a),f(b))$ on its abelianization, for
$f(T)=T+\epsilon(T^2-T)$. More generally, any such coordinatewise
lift with $f(1)=1$ forces $f$ to be multiplicative as well as
additive.
\end{statement}
\begin{proof}
The centre and commutator subgroup are the third-coordinate
$\mathbf G_a$. An automorphism induces some additive automorphism
$g$ on it. The commutator of $(a,0,0)$ and $(0,b,0)$ is
$(0,0,ab)$, independent of any chosen central parts of their lifts.
Compatibility with commutators therefore forces the polynomial
identity
\[
                    g(ab)=f(a)f(b).
\]
Setting $b=1$ gives $g=f$, and hence $f(ab)=f(a)f(b)$.
For $a=b=\omega$, however,
\[
 f(\omega^2)=\omega^2+\epsilon,
 \qquad f(\omega)^2=(\omega+\epsilon)^2=\omega^2.
\]
Their difference is the nonzero element $\epsilon$. This is a
finite ring counterexample, also checked by the executable above.

The same argument applies to a specified class-two quotient whose
commutator map contains this paired pair of free vector directions.
It does not identify that quotient with a canonical integral
relative completion. At the prime $2$, such an identification is
substantive: changing Hall coordinates in a nilpotent group can
introduce $\binom{T}{2}$, so a characteristic-zero exponential
description cannot simply be reduced modulo $2$.
\end{proof}

\begin{statement}{H.7}{Proved}{Action formulas and the joint full-prime abelian-stage kernel}
Let $\widehat C$ be the full profinite completion of $C$, and now
write $F,K_n$ for the closures of the preceding finite-index
subgroups. An actual continuous automorphism $\varphi$ preserving
$F,K_n$ acts on the full profinite abelian extension
$\widehat C/\overline{[K_n,K_n]}$ by
\[
 (v,h)\longmapsto(A_{\varphi,n}v+b_{\varphi,n}(h),
                         \theta_{\varphi,n}(h)),
\]
where $A_{\varphi,n}\in\mathrm{GL}_{n^2+1}(\widehat{\mathbf Z})$,
$b_{\varphi,n}:C/K_n\to\widehat{\mathbf Z}^{,n^2+1}$,
and $\theta_{\varphi,n}\in\operatorname{Aut}(C/K_n)$.
Its action is trivial exactly when $\theta=1$, $A=I$, and $b=0$.
Over the integral $\mathbf Z_\ell$ abelian-kernel model, the same
test uses the $\ell$-components of $A,b$ and the same finite
component permutation.

Furthermore
\[
 \bigcap_{n\geq1}\overline{[K_n,K_n]}
       =\overline{[F',F']}=F'',
\]
so the joint kernel on all full-prime abelian stages is precisely
the kernel of
$\operatorname{Aut}(\widehat C;F,(K_n))
\to\operatorname{Aut}(\widehat C/F'')$.
\end{statement}
\begin{proof}
The matrix is the continuous linear map on the free profinite
abelian group $K_n^{\mathrm{ab}}$. With word representatives
$s_h$, the translation is the abelianized class of
$\varphi(s_h)s_{\theta(h)}^{-1}$. Applying $\varphi$ to the
extension law gives the displayed affine formula and its
compatibilities. Equality with the identity is equivalent to the
three conditions, by inspecting components, their origins, and
their vector radicals. Tensoring the lattice with $\mathbf Z_\ell$
gives the asserted integral model and coefficient test.

Each $K_n$ contains $F'$, so $F''$ lies in the intersection.
For the reverse inclusion, replace $n$ by the cofinal sequence
$n!$, for which $\bigcap K_{n!}=F'$. Let $g\in F\setminus F''$.
There is a finite quotient $Q$ of $F/F''$ in which the image of
$g$ is nontrivial. By compactness the decreasing images of
$K_{n!}$ in $Q$ eventually equal the image of $F'$: a point in
every image has a compatible lift in the nested compact fibres.
This image is abelian because $Q$ is metabelian. Hence for all
sufficiently large $n$ the image of
$\overline{[K_{n!},K_{n!}]}$ is trivial, and it cannot contain
$g$. This proves the intersection formula. Finally an
automorphism is trivial at every stage exactly when
$\varphi(g)g^{-1}$ lies in that intersection for every
$g\in\widehat C$, which is the claimed joint-kernel equality.
\end{proof}

\begin{statement}{H.8}{Proved}{The nonlinear kernel direction is not an actual profinite signed action}
On an integral abelian-kernel stage every action induced by actual
profinite retained data is affine linear in the vector coordinates,
with coefficients in $\mathbf Z_\ell$. After base change to an
$\mathbf F_2$-algebra it remains affine linear. In particular the
nontrivial automorphism in H.5 is not induced by such data, even
though it fixes every discrete word point and extends through all
the abelian stages.
\end{statement}
\begin{proof}
This follows from H.7's formula. On the identity component the
translation is zero, so the induced vector-group map is linear.
The polynomial $T+\epsilon(T^2-T)$ has a nonzero quadratic
coefficient in $R[T]$, and therefore is not linear as a group-scheme
morphism. Checking only the embedded discrete word points would
miss this distinction, as H.5 proves. Thus the full ambient
polynomial lift kernel cannot be substituted for the image of
the retained profinite normalizer or conjugator source.
\end{proof}

\begin{statement}{H.9}{Proved}{The naive Hall model already loses the modular action at two}
Let $\mathscr H_5$ be the class-two group scheme over $\mathbf Z$
with coordinates $(v,z)\in\mathbf A^5\times\mathbf A^{10}$ and law
\[
 (v,z)(w,t)=\left(v+w,
       \bigl(z_{ij}+t_{ij}+v_iw_j\bigr)_{i<j}\right).
\]
Its integer points are a free class-two nilpotent lattice, with the
five coordinate generators corresponding to the Schreier basis of
H.4. There is no automorphism of its special fibre over $\mathbf F_2$
whose action on the designated quotient $\mathbf G_a^5$ is the
verified modular action of $\rho$ on $K^{\mathrm{ab}}$.
Thus this naive Hall model cannot be the radical of an integral
class-two modular extension realizing the specified lattice action.
\end{statement}
\begin{proof}
In characteristic two the square of every lift of $v$ is
\[
              (v,z)^2=(0,(v_iv_j)_{i<j}).
\]
The actual integral matrix in the q2 certificate is
\[
 A_\rho=\begin{pmatrix}
 0&-1&-1&1&0\\
 0&0&-1&0&0\\
 1&0&0&0&-1\\
 0&1&0&-1&-1\\
 0&0&-1&1&0
 \end{pmatrix}.
\]
It sends $e_2$ modulo two to $e_1+e_4$. The element $(e_2,0)$
has square one, whereas every point above $e_1+e_4$ has nonzero
square $e_1\wedge e_4$ in the centre. A homomorphism preserves
squares, giving a contradiction. The program
\texttt{code/fgh\_hall\_models.py} verifies the matrix against the
word representative for $\rho$ and checks all $2^{10}$ possible
central lifts over $\mathbf F_2$.
This is an obstruction to this explicitly specified integral model;
it is not an obstruction to the genuine profinite action, which
exists on the class-two quotient as a profinite group.
\end{proof}

\begin{statement}{H.10}{Proved}{The generic fibre and abelian stage do not determine the answer}
There are two smooth affine class-two group schemes over $\mathbf Z$
with isomorphic rational fibres and the same specified quotient
$\mathbf G_a^2$, for which the nonlinear additive automorphism of
H.5 has opposite extension behaviour after reduction modulo two.
Consequently an unspecified rational relative completion and its
integral first abelian stage do not determine H.3 at the prime two.
\end{statement}
\begin{proof}
For an integer $t$, define $\mathscr H_t=\mathbf A^3_{\mathbf Z}$ by
\[
 (a,b,c)(a',b',c')=(a+a',b+b',c+c'+t ab').
\]
The map $\mathscr H_{2,\mathbf Q}\to\mathscr H_{1,\mathbf Q}$,
$(a,b,c)\mapsto(a,b,c/2)$, is a group isomorphism compatible with
the quotients $(a,b,c)\mapsto(a,b)$. In characteristic two the
first model $\mathscr H_1$ is the Heisenberg group of H.6, so
$f(T)=T+\epsilon(T^2-T)$ cannot lift its coordinatewise action on
the quotient. The second model $\mathscr H_2$ becomes the additive
group $\mathbf G_a^3$, on which coordinatewise $f$ is an
automorphism. The program checks this latter identity over
$\mathbf F_4[\epsilon]/\epsilon^2$. Both claims concern the
specified quotient, without identifying it with the whole
abelianization of every special fibre.
\end{proof}

